\documentclass[10pt,a4paper]{amsart}
\usepackage[margin=19mm]{geometry}
\usepackage{amsmath,amssymb,mathtools}
\usepackage[scr=rsfs,cal=euler]{mathalfa}
\usepackage{xfrac}
\usepackage[unicode,hidelinks,bookmarks=false]{hyperref}
\newcommand{\ie}{i.e.\ }
\newcommand{\cf}{cf.\ }
\newcommand{\resp}{respectively\ }
\newcommand{\Subsection}{\subsection}
\providecommand{\nobreakdash}{\nobreak\hbox{-}}
\setlength{\emergencystretch}{2em}
\multlinegap0pt
\usepackage{xcolor}
\usepackage{amssymb}
\newtheorem{lemma}{Lemma}
\usepackage{cleveref}
\crefname{lemma}{Lemma}{Lemmas}
\newcommand{\F}{\mathbb{F}}
\title{Explicit Rank Extractors and Subspace Designs via Function Fields, with Applications to Strong Blocking Sets}
\author{Zeyu Guo, Roshan Raj, Chong Shangguan, Zihan Zhang}
\date{Source: arXiv:2604.13431v2 (CC BY 4.0)}
\begin{document}
\maketitle

\noindent\emph{Source and licence.} Explicit Rank Extractors and Subspace Designs via Function Fields, with Applications to Strong Blocking Sets. Version arXiv:2604.13431v2 (CC BY 4.0), distributed by arXiv under the Creative Commons Attribution 4.0 International licence, http://creativecommons.org/licenses/by/4.0/, as recorded in the arXiv licence metadata for this exact version and shown on the abstract page https://arxiv.org/abs/2604.13431v2. Full licence notice: \texttt{licenses/arxiv-2604.13431-CC-BY-4.0.txt}.\par\smallskip

\noindent\textit{Editorial scope.} Counted unit: the article's lemma lem:FS-deg (source0.tex lines 1196--1203). The article's complete original proof of it is delivered with it (source0.tex lines 1204--1392, one \texttt{proof} environment of 189 lines). No mathematics is written, repaired or re-derived: the statement and the proof are the article's own lines, in the article's own notation, and no retained line is substituted or re-flowed.  Retained uncounted as the source-local context the proof needs: lines 542--556; lines 708--710; lines 1167--1169; lines 1171--1173.  Nothing was omitted from any retained span, so the package carries no omission record.  No retained line contains a citation, so the wrapper carries no bibliography and no bibliography entry.  No retained line contains a figure environment, an image-inclusion command or drawing code, so no image is delivered and the package declares no asset dependency.  Editorial formatting only, in addition: an AMS \texttt{amsart} preamble that restates the article's own declarations and macros used by the retained lines, namely the environments \texttt{lemma} on separate counters, together with the standard packages those lines need.  The retained environments print in this document as Lemma 1 in order of appearance, whereas the article prints the same items with the numbering of its own full document.  The complete native source of the article as distributed in the arXiv e-print for this exact version is bundled unchanged as \texttt{sources/198473/source0.tex}.
\begin{lemma}\label{lem:unique}
Let $M \in \F^{r \times k}$ be a full-rank matrix over a field $\F$, with $r \le k$. 
Let $I \subseteq [k]$ be the subset produced by the following procedure.
Initialize $I = \emptyset$. For $i = k, k-1, \dots, 1$, add $i$ to $I$ if the $i$-th column of $M$ does not lie in the $\F$-linear span of the columns indexed by the current set $I$.

Then $|I| = r$ and $M_I$ is nonsingular. Moreover, for every $J \subseteq [k]$ of size $r$ such that $M_J$ is nonsingular, the following hold:
\begin{enumerate}
\item $i_\ell \ge j_\ell$ for all $\ell \in [r]$, where $i_\ell$ and $j_\ell$ denote the $\ell$-th smallest elements of $I$ and $J$, respectively.
\item Let $w_1 > w_2 > \cdots > w_k$ be integers. Then
\[
\sum_{i \in I} w_i \le \sum_{i \in J} w_i,
\]
with equality if and only if $I = J$.
\end{enumerate}
\end{lemma}

\begin{lemma}\label{lem:triangle}
If $v(x)\neq v(y)$, then $v(x+y)=\min\{v(x),v(y)\}$.
\end{lemma}

\begin{equation}\label{eq:dec-ord-1}
0\geq v_{P_\infty}(f_1) > \cdots > v_{P_\infty}(f_k) \geq -d_1
\end{equation}

\begin{equation}\label{eq:dec-ord-2}
0\geq v_{P_\infty}(g_1) > \cdots > v_{P_\infty}(g_h) \geq -d_2.
\end{equation}

\begin{lemma}\label{lem:FS-deg}
Let $M\in \mathbb{F}_q^{k\times r}$ be full rank. Then $\det(EM)\in F$ is nonzero, $v_P(\det(EM))\geq 0$ for all places $P\neq P_\infty$, and 
\[
v_{P_\infty}(\det(EM))\geq -r\left(
k-r+\frac{r-1}{2}\left\lceil \frac{k}{q-1}\right\rceil+\frac{r+1}{2}g 
\right).
\]
\end{lemma}

\begin{proof}
For any place $P\neq P_\infty$, we have $v_P(\det(EM))\geq 0$ since $v_P(a)\geq 0$ whenever $a$ is an entry of $E$ or an entry of $M$.

By the Cauchy--Binet formula,
\[
\det(EM)=\sum_{I\subseteq [k],\, |I|=r}\det(E_I)\det(M^I).
\]

Fix $I=\{j_1<\cdots<j_r\}\subseteq [k]$. The matrix $E_I$ has the form
\[
E_I=\bigl((\gamma^{\beta(j_t)} g_{\alpha(j_t)})^{i-1} f_{j_t}\bigr)_{i,t\in [r]}.
\]
Factoring out $f_{j_t}$ from each column and using multilinearity of the determinant, we obtain
\[
\det(E_I)=\left(\prod_{u=1}^r f_{j_u}\right)
\det\bigl((\gamma^{\beta(j_t)} g_{\alpha(j_t)})^{i-1}\bigr)_{i,t\in [r]}.
\]
The remaining determinant is a Vandermonde determinant, and hence
\begin{equation}\label{eq:expansion-EI}
\det(E_I)=\left(\prod_{u=1}^r f_{j_u}\right)
\prod_{1\le t_1<t_2\le r}
\left(\gamma^{\beta(j_{t_2})} g_{\alpha(j_{t_2})}-\gamma^{\beta(j_{t_1})} g_{\alpha(j_{t_1})}\right).
\end{equation}

\medskip
We now analyze the valuation of each difference term. For $t_1<t_2$, consider
\[
\gamma^{\beta(j_{t_2})} g_{\alpha(j_{t_2})}
-
\gamma^{\beta(j_{t_1})} g_{\alpha(j_{t_1})}.
\]
There are two cases.

\smallskip
\noindent\textbf{Case 1:} $\alpha(j_{t_1})=\alpha(j_{t_2})$.
Then $\beta(j_{t_1})<\beta(j_{t_2})$, and
\[
v_{P_\infty}\!\left(
\gamma^{\beta(j_{t_2})} g_{\alpha(j_{t_2})}
-
\gamma^{\beta(j_{t_1})} g_{\alpha(j_{t_1})}
\right)
=
v_{P_\infty}\!\left(
(\gamma^{\beta(j_{t_2})}-\gamma^{\beta(j_{t_1})}) g_{\alpha(j_{t_2})}
\right)
=
v_{P_\infty}(g_{\alpha(j_{t_2})}).
\]

\smallskip
\noindent\textbf{Case 2:} $\alpha(j_{t_1})<\alpha(j_{t_2})$.
Then by \eqref{eq:dec-ord-2},
\[
v_{P_\infty}(g_{\alpha(j_{t_2})})<v_{P_\infty}(g_{\alpha(j_{t_1})}),
\]
and hence, by \cref{lem:triangle},  $v_{P_\infty}\!\left(
\gamma^{\beta(j_{t_2})} g_{\alpha(j_{t_2})}
-
\gamma^{\beta(j_{t_1})} g_{\alpha(j_{t_1})}
\right)
=
v_{P_\infty}(g_{\alpha(j_{t_2})})$.

\smallskip
Thus in both cases, 
\begin{equation}\label{eq:diff}
v_{P_\infty}\!\left(
\gamma^{\beta(j_{t_2})} g_{\alpha(j_{t_2})}
-
\gamma^{\beta(j_{t_1})} g_{\alpha(j_{t_1})}
\right)
=
v_{P_\infty}(g_{\alpha(j_{t_2})}).
\end{equation}

\medskip
For each $I=\{j_1<\cdots<j_r\}$, define
\[
T_I=\det(E_I)\det(M^I).
\]
From \eqref{eq:expansion-EI} and the definition of $T_I$, we obtain
\[
v_{P_\infty}(T_I)
=
\sum_{u=1}^r v_{P_\infty}(f_{j_u})
+
\sum_{1\le t_1<t_2\le r}
v_{P_\infty}\!\left(
\gamma^{\beta(j_{t_2})} g_{\alpha(j_{t_2})}
-
\gamma^{\beta(j_{t_1})} g_{\alpha(j_{t_1})}
\right)
+
v_{P_\infty}(\det(M^I)).
\]
By \eqref{eq:diff}, each term in the second sum equals
$v_{P_\infty}(g_{\alpha(j_{t_2})})$, and hence
\begin{equation}\label{eq:expansion-TI}
\begin{aligned}
v_{P_\infty}(T_I)
&=
\sum_{u=1}^r v_{P_\infty}(f_{j_u})
+
\sum_{t=1}^r (t-1)v_{P_\infty}(g_{\alpha(j_t)})
+
v_{P_\infty}(\det(M^I)).\\
&=\begin{cases}
\sum_{u=1}^r v_{P_\infty}(f_{j_u})
+
\sum_{t=1}^r (t-1)v_{P_\infty}(g_{\alpha(j_t)}) & \text{if } \det(M^I)\neq 0,\\
\infty & \text{otherwise.}
\end{cases}
\end{aligned}
\end{equation}

Let $I^*\subseteq [k]$ be the subset produced by the following greedy procedure:
initialize $I^*=\emptyset$, and for $i=k,k-1,\dots,1$, add $i$ to $I^*$ if the $i$-th column of $M$ does not lie in the $\F$-linear span of the columns indexed by the current set $I^*$.
By \cref{lem:unique}, we have $|I^*|=r$ and $\det(M^{I^*})\neq 0$.

First, we show that $I^*$ uniquely minimizes
the first sum $\sum_{u=1}^r v_{P_\infty}(f_{j_u})$ in \eqref{eq:expansion-TI}
among all subsets $I\subseteq [k]$ of size $r$ satisfying $\det(M^I)\neq 0$.
Indeed, by \eqref{eq:dec-ord-1}, the sequence
\[
v_{P_\infty}(f_1) > v_{P_\infty}(f_2) > \cdots > v_{P_\infty}(f_k)
\]
is strictly decreasing. Applying the second item of \cref{lem:unique} with weights $w_i = v_{P_\infty}(f_i)$, we conclude that
\[
\sum_{i\in I^*} v_{P_\infty}(f_i)
\le \sum_{i\in J} v_{P_\infty}(f_i)
\]
for every subset $J\subseteq [k]$ of size $r$ such that $\det(M^J)\neq 0$, with equality if and only if $J=I^*$. This proves the claim.
 

Next, we show that $I^*$ also minimizes the second sum $\sum_{t=1}^r (t-1)v_{P_\infty}(g_{\alpha(j_t)})$
among all subsets $I=\{j_1<\cdots<j_r\}$ satisfying $\det(M^I)\neq 0$.

Let $I=\{j_1<\cdots<j_r\}$ be any such subset, and write $I^*=\{j_1^*<\cdots<j_r^*\}$.
By the first item of \cref{lem:unique}, we have $j_t^*\ge j_t$ for all $t\in [r]$.
Since $\alpha(\cdot)$ is nondecreasing and the sequence
$v_{P_\infty}(g_1) > \cdots > v_{P_\infty}(g_h)$
is strictly decreasing by \eqref{eq:dec-ord-2}, it follows that
\[
v_{P_\infty}(g_{\alpha(j_t^*)})
\le
v_{P_\infty}(g_{\alpha(j_t)})
\qquad\text{for all } t\in [r].
\]
Multiplying by $(t-1)\ge 0$ and summing over $t$, we obtain
\[
\sum_{t=1}^r (t-1)v_{P_\infty}(g_{\alpha(j_t^*)})
\le
\sum_{t=1}^r (t-1)v_{P_\infty}(g_{\alpha(j_t)}),
\]
as claimed.

Combining the two parts above, we conclude that $I^*$ uniquely minimizes $v_{P_\infty}(T_I)$ among all subsets $I\subseteq [k]$ of size $r$.

Returning to the Cauchy--Binet expansion
\[
\det(EM)=\sum_{I\subseteq [k],\, |I|=r} T_I,
\]
we have $\det(M^{I^*})\neq 0$, and hence \eqref{eq:expansion-TI} implies that
$v_{P_\infty}(T_{I^*}) < \infty$.
Since $T_{I^*}$ is the unique term attaining the minimum valuation, by \cref{lem:triangle},
\[
v_{P_\infty}(\det(EM)) = v_{P_\infty}(T_{I^*}) < \infty,
\]
and thus $\det(EM)\neq 0$.

Moreover, writing $I^*=\{j_1^*<\cdots<j_r^*\}$ and using \eqref{eq:expansion-TI}, we have
\begin{align*}
v_{P_\infty}(\det(EM))
&=
\sum_{u=1}^r v_{P_\infty}(f_{j_u^*})
+
\sum_{t=1}^r (t-1)v_{P_\infty}(g_{\alpha(j_t^*)})\\
&\stackrel{\eqref{eq:dec-ord-1},\eqref{eq:dec-ord-2}}{\ge}
-\sum_{u=1}^r(d_1-r+u)-\sum_{t=1}^r (t-1)d_2\\
&=-\sum_{u=1}^r(k-1+g-r+u)-\sum_{t=1}^r (t-1)\left(\left\lceil\frac{k}{q-1}\right\rceil-1+g\right),
\end{align*}
which simplifies to
\[
v_{P_\infty}(\det(EM))\geq -r\left(
k-r+\frac{r-1}{2}\left\lceil \frac{k}{q-1}\right\rceil+\frac{r+1}{2}g 
\right).\qedhere
\]
\end{proof}

\end{document}
